% from main.tex
%%%%%%%%%%  scan idx 101  |  folio 88  %%%%%%%%%%
sens ordinaire. La sous-catégorie pleine \(\mathcal C_1\) de \(\mathcal C\) formée des \(A\)-modules de type fini est quasi-relevable dans \(\mathcal C\) (mais non relevable en général)~; si \(A\) est noethérien, \(\mathcal C_1\) est stable par noyau d'épimorphisme. La sous-catégorie pleine \(\mathcal C'_0\) de \(\mathcal C\) formée des \(A\)-modules libres de type fini est relevable dans \(\mathcal C\) et quasi-stable (non stable en général) par noyau d'épimorphisme.

\textbf{b)} Soient \((S,\mathcal O_S)\) un espace annelé en anneaux locaux, \(S\) le site des ouverts de \(S\), \(\mathcal C\) la \(S\)-catégorie fibrée (et même scindée (SGA 1 VI 9)) définie par le foncteur \(U\longmapsto\) (catégorie des \(\mathcal O_U\)-Modules à gauche)~; \(\mathcal C\) est une \(S\)-catégorie abélienne plate (1.1). Soit \(\mathcal C_0\) la sous-\(S\)-catégorie additive strictement pleine de \(\mathcal C\) définie par \(U\longmapsto\) (catégorie des \(\mathcal O_U\)-Modules libres de type fini). Un \(\mathcal O_S\)-Module \(M\) est localement de \(\mathcal C_0\)-type fini si et seulement si \(M\) est de type fini au sens ordinaire (EGA \(0_{\mathrm I}\) 5.2). De (EGA \(0_{\mathrm I}\) 5.2.7 et 5.4.9) résulte que \(\mathcal C_0\) est localement stable par noyau d'épimorphisme et localement relevable dans \(\mathcal C\).

\textbf{c)} Soient \(G\) un groupe, \(S\) le topos des \(G\)-ensembles (SGA 4 V 3.11) \(k\) un corps (considéré comme Anneau constant de \(S\)), \(\mathcal C\) la \(S\)-catégorie scindée définie par le foncteur \(U\longmapsto\) (catégorie des \(k_U\)-Modules à gauche)~; \(\mathcal C\) est une \(S\)-catégorie abélienne plate. Soit \(\mathcal C_0\) la sous-\(S\)-catégorie additive strictement pleine de \(\mathcal C\) définie par \(U\longmapsto\) (catégorie des \(k_U\)-Modules libres de type fini).
Comme la localisation dans \(S\) équivaut à l'oubli des opérations de \(G\), il est clair qu'un \(G\)-\(k\)-module \(M\) est localement de \(\mathcal C_0\)-type fini si et seulement si \(M\) est de dimension finie comme espace vectoriel sur \(k\), et que \(\mathcal C_0\) est localement stable par noyau d'épimorphisme et localement relevable dans \(\mathcal C\).

On laisse au lecteur le soin de regarder la variante «~galoisienne~» de c) où \(G\) est un groupe profini et \(S\) le topos des ensembles où \(G\) opère continûment.

